\chapter{Project Requirements}
\label{sec:requirements}

This chapter provides a comprehensive specification of the functional and non-functional requirements for the personalized-ai-storybook system as implemented in FYP-2. The requirements reflect the complete system including all FYP-2 additions: TTS, STT, styled PDF export, personalized and simple story enhancement, evaluation metrics, and UI improvements.

\section{Use-case Analysis}

The personalized-ai-storybook platform supports comprehensive operational workflows encompassing story generation, audio narration, voice input, personalized character rendering, interactive chat, idea workshop, and quality evaluation.

\subsection{Use-Case: Generate Story from Text Prompt (UC1)}
\textbf{Primary Actor}: Registered User \\
\textbf{Precondition}: User is authenticated and on the story generation page. \\
\textbf{Main Flow}:
\begin{enumerate}
    \item User selects ``Simple'' mode, enters a text prompt, and chooses genre/style/length.
    \item System validates and enhances the prompt via Prompt Agent.
    \item Director Agent orchestrates Writer → Reviewer → Editor pipeline.
    \item System generates illustrated scenes and exports a styled PDF.
    \item Story is displayed in the open-book UI with audio controls.
\end{enumerate}
\textbf{Postcondition}: Story is stored in the database; PDF is saved in \texttt{generated/exports/}.

\subsection{Use-Case: Generate Story from Voice Input (UC2)}
\textbf{Primary Actor}: Registered User \\
\textbf{Precondition}: User is on the story generation page; microphone is accessible. \\
\textbf{Main Flow}:
\begin{enumerate}
    \item User clicks the voice input button and speaks the story prompt.
    \item System records audio, sends WAV to \texttt{/api/stt} endpoint.
    \item Whisper STT transcribes speech to text.
    \item Transcribed text is pre-filled in the prompt field.
    \item User confirms and story generation proceeds as in UC1.
\end{enumerate}

\subsection{Use-Case: Listen to Story Narration (UC3)}
\textbf{Primary Actor}: Registered User \\
\textbf{Precondition}: Story has been generated and displayed. \\
\textbf{Main Flow}:
\begin{enumerate}
    \item User selects voice (male/female) from the audio control panel.
    \item User clicks ``Play'' for the current page.
    \item System calls \texttt{/api/tts} to generate/retrieve cached WAV audio.
    \item Audio plays through the browser; user can pause or switch pages.
\end{enumerate}

\subsection{Use-Case: Generate Personalized Story (UC4)}
\textbf{Primary Actor}: Registered User \\
\textbf{Precondition}: User selects ``Personalized'' mode and uploads a reference photo. \\
\textbf{Main Flow}:
\begin{enumerate}
    \item User uploads a photo in the personalized mode panel.
    \item System processes photo through IP-Adapter FaceID Plus v2.
    \item Story is generated; all illustrations maintain facial likeness using WebUI API.
    \item Personalized story is displayed and exported.
\end{enumerate}

\subsection{Use-Case: View Evaluation Metrics (UC5)}
\textbf{Primary Actor}: Registered User \\
\textbf{Precondition}: Story has been generated; evaluation has run in the background. \\
\textbf{Main Flow}:
\begin{enumerate}
    \item Evaluation Agent runs asynchronously after story generation.
    \item Computes: image-text alignment, visual consistency, text coherence, readability, story structure.
    \item Results stored in \texttt{generated/evaluations/} as JSON.
    \item User can view quality metrics in the story details panel.
\end{enumerate}

\subsection{Use-Case: Chat with Story Character (UC6)}
\textbf{Primary Actor}: Registered User \\
\textbf{Main Flow}: User selects a character from the story; types a message; Chatbot Agent responds in character using story context.

\subsection{Use-Case: Idea Workshop (UC7)}
\textbf{Primary Actor}: Registered User \\
\textbf{Main Flow}: User starts a workshop session; the LangChain-based Idea Workshop Agent collects story requirements through multi-turn conversation; confirms all fields; generates the story.

\subsection{Use-Case: Export Story to PDF (UC8)}
\textbf{Primary Actor}: Registered User \\
\textbf{Main Flow}: Upon story completion, system automatically calls the EditorAgent PDF export; the styled PDF with leather cover and notebook-ruled text pages is saved and available for download.


\section{Functional Requirements}

\subsection{Multi-Agent Story Generation Module}
\begin{enumerate}
    \item \textbf{FR1.1}: The system shall accept user prompts ranging from 5 to 500 characters for story generation.
    \item \textbf{FR1.2}: The Prompt Agent shall validate and classify prompts as normal, detailed, invalid, or nonsense.
    \item \textbf{FR1.3}: The Prompt Agent shall enhance simple prompts by adding contextual details.
    \item \textbf{FR1.4}: The Director Agent shall orchestrate the sequential execution of Prompt, Writer, Reviewer, and Editor agents.
    \item \textbf{FR1.5}: The Writer Agent shall generate structured stories using Ollama LLM (llama3.1:8b-instruct-q4\_K\_M) with title, setting, characters, and exactly 3 scenes.
    \item \textbf{FR1.6}: The Writer Agent shall retry up to 2 times if JSON parsing fails.
    \item \textbf{FR1.7}: The Reviewer Agent shall validate content for coherence, age-appropriateness, and educational value.
    \item \textbf{FR1.8}: The Editor Agent shall finalize the story and trigger PDF export.
    \item \textbf{FR1.9}: The Evaluation Agent shall run asynchronously in the background after generation.
    \item \textbf{FR1.10}: The system shall output stories in comprehensive JSON format with metadata and scene image paths.
    \item \textbf{FR1.11}: The Content Safety module shall block explicit, violent, or inappropriate prompts and image descriptions.
    \item \textbf{FR1.12}: The Idea Workshop Agent shall collect story requirements through multi-turn LangChain conversation before generating.
\end{enumerate}

\subsection{Image Generation Module}
\begin{enumerate}
    \item \textbf{FR2.1}: The Image Agent shall generate illustrations for each scene using Stable Diffusion 1.5 with Realistic Vision checkpoint.
    \item \textbf{FR2.2}: The system shall employ a two-pass generation: txt2img at 512$\times$512 followed by img2img at 768$\times$768.
    \item \textbf{FR2.3}: The system shall use DPMSolver++ 2M Karras scheduler for optimized sampling.
    \item \textbf{FR2.4}: In Personalized mode, the system shall apply IP-Adapter FaceID Plus v2 embeddings via WebUI API.
    \item \textbf{FR2.5}: The system shall pause/resume Ollama service during image generation to manage GPU memory.
    \item \textbf{FR2.6}: The system shall continue story generation with text only if image generation fails.
\end{enumerate}

\subsection{Text-to-Speech Module}
\begin{enumerate}
    \item \textbf{FR3.1}: The system shall generate WAV audio narration for each story scene using Piper TTS (offline ONNX models).
    \item \textbf{FR3.2}: The system shall support male (en\_US-ryan-medium) and female (en\_US-lessac-medium) voice options.
    \item \textbf{FR3.3}: The system shall cache generated audio files to avoid redundant synthesis on re-request.
    \item \textbf{FR3.4}: TTS audio shall be served via static file routes accessible to the frontend.
    \item \textbf{FR3.5}: The system shall expose a \texttt{POST /api/tts} endpoint accepting story\_id, scene\_number, and voice parameters.
\end{enumerate}

\subsection{Speech-to-Text Module}
\begin{enumerate}
    \item \textbf{FR4.1}: The system shall transcribe WAV audio files to text using OpenAI Whisper (base model, offline).
    \item \textbf{FR4.2}: The system shall read WAV files using Python's built-in \texttt{wave} module without requiring ffmpeg.
    \item \textbf{FR4.3}: The system shall resample audio to 16kHz using numpy linear interpolation for Whisper compatibility.
    \item \textbf{FR4.4}: The system shall convert stereo recordings to mono automatically.
    \item \textbf{FR4.5}: The system shall expose a \texttt{POST /api/stt} endpoint accepting an audio file upload.
    \item \textbf{FR4.6}: The system shall return the transcribed text and detected language in the response.
\end{enumerate}

\subsection{PDF Rendering Module}
\begin{enumerate}
    \item \textbf{FR5.1}: The system shall generate a styled PDF storybook matching the open-book UI design.
    \item \textbf{FR5.2}: The PDF shall include a cover page with leather/cream layout, story title in gold typography, and last scene image.
    \item \textbf{FR5.3}: Each scene shall be rendered as a spread: illustration on the left with white photo mat; text on the right with red margin line and notebook rulings.
    \item \textbf{FR5.4}: Text shall auto-scale (14pt down to 10pt) to prevent overflow.
    \item \textbf{FR5.5}: The PDF shall use Georgia font with fallback to Times-Roman.
    \item \textbf{FR5.6}: The system shall use A4 landscape orientation with a leather-brown background frame.
    \item \textbf{FR5.7}: Each scene page shall include a page number pill indicator.
\end{enumerate}

\subsection{Evaluation Module}
\begin{enumerate}
    \item \textbf{FR6.1}: The Evaluation Agent shall assess image-text alignment using CLIP cosine similarity (OpenAI CLIP ViT-B/32).
    \item \textbf{FR6.2}: The system shall evaluate visual consistency across scene images using CLIP image embeddings.
    \item \textbf{FR6.3}: The system shall evaluate text coherence between consecutive scenes using Sentence Transformers.
    \item \textbf{FR6.4}: The system shall assess readability using Flesch-Kincaid grade level (target: grades 3--6 for ages 7--10).
    \item \textbf{FR6.5}: The system shall validate story structure (beginning, climax, ending) using a heuristic scorer.
    \item \textbf{FR6.6}: In Personalized mode, the system shall evaluate character face consistency using InsightFace embeddings.
    \item \textbf{FR6.7}: The system shall compute a weighted overall score (0--100) and save results as JSON.
\end{enumerate}

\subsection{User Interface and Frontend Module}
\begin{enumerate}
    \item \textbf{FR7.1}: The frontend shall provide a Next.js-based responsive open-book UI.
    \item \textbf{FR7.2}: The story display shall present a left page (illustration) and right page (story text) with page-flip navigation.
    \item \textbf{FR7.3}: Each page shall include audio controls (play/pause button, voice selector).
    \item \textbf{FR7.4}: The UI shall include a voice input button for speech-to-text story prompt submission.
    \item \textbf{FR7.5}: The system shall provide theme toggle for light/dark mode.
    \item \textbf{FR7.6}: The UI shall support genre, style, and length customization dropdowns.
    \item \textbf{FR7.7}: The UI shall display custom loading animations during story generation.
    \item \textbf{FR7.8}: Simple and Personalized mode tabs shall be clearly separated in the UI.
\end{enumerate}


\section{Non-Functional Requirements}

\subsection{Performance}
\begin{itemize}
    \item \textbf{PER-1}: Story text generation shall complete within 10--15 seconds on systems with 16GB RAM.
    \item \textbf{PER-2}: Image generation per scene shall complete within 35--45 seconds on GPU (NVIDIA GTX 1060 or equivalent).
    \item \textbf{PER-3}: TTS audio generation per scene shall complete within 2--5 seconds using Piper.
    \item \textbf{PER-4}: STT transcription shall complete within 3--8 seconds for a 15-second voice prompt.
    \item \textbf{PER-5}: Evaluation Agent shall run asynchronously and not block the story generation response.
    \item \textbf{PER-6}: PDF export shall complete within 3--5 seconds for a 3-scene story.
\end{itemize}

\subsection{Reliability}
\begin{itemize}
    \item \textbf{REL-1}: The Writer Agent shall successfully generate valid JSON for 90\% of reasonable prompts.
    \item \textbf{REL-2}: The system shall retry up to 2 times on JSON parsing failures.
    \item \textbf{REL-3}: The system shall deliver text-only stories when image generation fails (graceful degradation).
    \item \textbf{REL-4}: TTS shall return an empty string rather than crashing if Piper synthesis fails.
    \item \textbf{REL-5}: STT shall raise a descriptive exception rather than silently failing.
    \item \textbf{REL-6}: The Evaluation Agent shall handle missing images gracefully with default scores.
\end{itemize}

\subsection{Usability}
\begin{itemize}
    \item \textbf{USE-1}: Users with no technical background shall be able to generate a story within 3 interactions.
    \item \textbf{USE-2}: The UI shall provide clear visual feedback (loading spinner, progress messages) during all operations.
    \item \textbf{USE-3}: The system shall display user-friendly error messages when generation fails.
    \item \textbf{USE-4}: Voice input shall be accessible via a single button click without configuration.
    \item \textbf{USE-5}: Audio narration shall start within 1 second of the user clicking play (from cache).
\end{itemize}

\subsection{Security}
\begin{itemize}
    \item \textbf{SEC-1}: All AI processing shall occur locally without transmitting user data to external servers.
    \item \textbf{SEC-2}: User authentication shall use JWT tokens with bcrypt password hashing.
    \item \textbf{SEC-3}: User-uploaded reference photos shall not be stored permanently without user consent.
    \item \textbf{SEC-4}: The Content Safety module shall block inappropriate content before generation.
\end{itemize}

\subsection{Maintainability}
\begin{itemize}
    \item \textbf{MNT-1}: Each agent module shall be independently implementable and testable.
    \item \textbf{MNT-2}: TTS and STT managers shall use a singleton pattern for resource efficiency.
    \item \textbf{MNT-3}: All evaluation scores shall be in a unified 0--100 integer scale for easy interpretation.
    \item \textbf{MNT-4}: The codebase shall include inline comments and docstrings for all key functions.
\end{itemize}

\subsection{Portability}
\begin{itemize}
    \item \textbf{PORT-1}: The system shall operate on Windows 10/11 systems (primary development platform).
    \item \textbf{PORT-2}: The STT module shall not require ffmpeg installation (pure Python WAV reader).
    \item \textbf{PORT-3}: The TTS module shall use offline ONNX Piper models, requiring no internet.
    \item \textbf{PORT-4}: The system shall function on both CPU-only and GPU-enabled hardware.
\end{itemize}
